\scanpage{24}
2. Si on se donne une section $s$ de l'enveloppe affine, i.e. un splitting de (2), i.e.
un iso de $\underline{O}_S$-Modules
\[
B^{\deleted{\unsure{\vee}}}\simeq\underline{O}_S\oplus\check L,\tag*{(11)}
\]
alors la structure de $\underline{O}_S$-Algèbre de $B$ est déterminée
par la donnée de la multiplication
\[
\check L\otimes\check L\longrightarrow B=\underline{O}_S+\check L
\]
qui sera donnée elle-\unsure{=|même} par deux hom
\[
\left\{\begin{aligned}
b&\in\Gamma\underline{\mathrm{Hom}}(L^{\otimes(-2)},L^{-1})\simeq\Gamma L\\
c&\in\Gamma\underline{\mathrm{Hom}}(L^{\otimes(-2)},\underline{O}_S)\simeq\Gamma L^{\otimes2}.
\end{aligned}\right. \tag*{(12)}
\]
De façon précise, on aura
\[
\left\{\begin{aligned}
(\lambda\oplus s(u))(\mu\oplus s(v))&=(\lambda\mu\inserted{\deleted{\unsure{+}}-cuv})\oplus s(\lambda v+\mu u\deleted{\unsure{-}}+buv).\\
\text{i.e. }s(u)s(v)&=-cuv\oplus s(buv)\qquad\text{i.e. }\deleted{\unsure{uv=}}
\end{aligned}\right. \tag*{(13)}
\]
\margin{attention à la convention de signe !}
\deleted{On aura} Calculons $\operatorname{tr}_{\Omega/S}(\lambda+u)$ et $N_{\Omega/S}(\lambda+u)$.

\deleted{(14)} On aura
\[
\operatorname{tr}(\lambda+u)=2\lambda+\operatorname{tr}u
\]
Or la multiplication $\mu_u$ par $u$ est donné par
\[
\mu\oplus v\longmapsto\inserted{-}cuv\oplus s(\mu u+buv)
\]
\margin{signe $-$ !}
i.e. par la matrice
\[
\mu_u=\begin{pmatrix}0&-cu\\u&-bu\end{pmatrix},\quad
\mu_{\lambda\oplus u}=\begin{pmatrix}\lambda&-cu\\u&\lambda+bu\end{pmatrix}
\]
donc
\[
\left\{\begin{aligned}
\operatorname{tr}u&=\deleted{\unsure{-}}bu\\
Nu&=\deleted{\unsure{-}}cu^2
\end{aligned}\right. \tag*{(14)}
\]
d'où
\[
\left\{\begin{aligned}
\operatorname{tr}(\lambda\oplus u)&=2\lambda\deleted{\unsure{-}}+bu\\
N(\lambda\oplus u)&=\lambda^2+\lambda bu\deleted{\unsure{-}}+cu^2
\end{aligned}\right. \tag*{(15)}
\]
\margin{En termes de $s$ définissant un iso $F\simeq\check{\mathbb V}(L)$, on a un plongmt $\Omega\subset\check{\mathbb V}(L)$, et $\Omega$ est caractérisé dans $\check{\mathbb V}(L)$ comme le ss schéma des sections $x$ \unsure{tels|telles} que $\boxed{x^2-bx+c=0}$.}
% margin note written sideways
\margin{La cat. des $\Omega$ (ou $B$) \deleted{\unsure{munis}} munis d'un scindage équivaut à celle des \unsure{triples|triplets} $(L,b,c)$ \ldots .}

Voyons la dépendance de $b=b(s)$ et $c=c(s)$ \unsure{p.r.|par r.} à $s$. Soit \deleted{\unsure{en}}
\[
s'=s+\tau\qquad \tau\in\Gamma(E)\qquad\text{i.e. }s'(u)=s(u)+\tau u
\quad\text{pour }u\in\Gamma\check L \tag*{(16)}
\]

% archive sheet number: 23

